\documentclass[../main.tex]{subfiles}

\begin{document}

\section{Glossary}
\label{app:glossary}

This appendix explains specialized vocabulary used in the report for readers
who do not work in security, privacy, software assurance, or HDF5 development.
Definitions describe how a term is used \emph{in this report}; they do not add
a finding, change a rating, or replace the controlling scope and decision rules
in Sections~\ref{sec:assessment-scope}--\ref{sec:risk-rubric}. Product and
project names are generally omitted unless the name itself is used as
technical shorthand.

\subsection{Abbreviations}

\begingroup
\small
\setlength{\tabcolsep}{3pt}
\renewcommand{\arraystretch}{1.12}
\begin{longtable}{@{}>{\raggedright\arraybackslash}p{0.15\linewidth}>{\raggedright\arraybackslash}p{0.25\linewidth}p{0.54\linewidth}@{}}
\hline
\textbf{Short form} & \textbf{Expansion} & \textbf{Meaning in this report} \\
\hline
\endfirsthead
\hline
\textbf{Short form} & \textbf{Expansion} & \textbf{Meaning in this report} \\
\hline
\endhead
\hline
\endfoot

ABI & Application Binary Interface &
The low-level contract by which compiled software calls a library, including
calling conventions, data layout, and exported symbols. Two builds can expose
similar source APIs but still be ABI-incompatible. \\

AFL & American Fuzzy Lop &
A coverage-guided fuzzing tool family. The build appendix mentions AFL support
as one available test configuration, not as evidence that every parser path was
fuzzed. \\

AI/ML & Artificial Intelligence / Machine Learning &
Application domains discussed here chiefly because model loaders and model
artifacts can cross an HDF5 trust boundary. \\

API & Application Programming Interface &
The functions, types, and rules through which software uses a library or
service. The HDF5 C API is distinct from the HDF5 on-disk format. \\

ASan & AddressSanitizer &
A runtime testing tool that detects many invalid memory accesses, including
some buffer overflows and uses of freed memory. \\

CI / CI/CD & Continuous Integration / Continuous Integration and Delivery &
Automated building and testing of changes; CD adds automated packaging or
delivery stages. Passing CI is evidence only for the configurations and tests
that actually ran. \\

CLI & Command-Line Interface &
A program operated from a shell, such as \code{h5dump} or \code{h5repack}. \\

CPU & Central Processing Unit &
The processor resource that can be exhausted by excessive or non-terminating
work. \\

CVE & Common Vulnerabilities and Exposures &
A public identifier for a reported vulnerability record. A CVE identifier does
not by itself establish root cause, affected HDF5 versions, reachability in a
particular application, or this report's risk rating. \\

CVSS & Common Vulnerability Scoring System &
A standardized vulnerability-severity scoring method. This report does not
convert CVSS scores into its impact/exposure ratings. \\

CWE & Common Weakness Enumeration &
A catalog of software weakness classes. For example, \code{CWE-20} denotes
improper input validation; a CWE describes a class, not a particular affected
release. \\

CWD & Current Working Directory &
The process directory used to resolve many relative paths. Implicit CWD lookup
can make an external-file or plugin search depend on who launched the process
and from where. \\

DNS & Domain Name System &
The network service that resolves names to addresses. DNS queries and logs can
expose which external data services a workflow contacts. \\

DoS & Denial of Service &
Loss of availability caused by a crash, hang, excessive work, exhausted
resources, or another condition that prevents useful service. \\

ENOSPC & No space left on device &
The operating-system error reported when a storage target cannot accept more
data. The report also uses it as shorthand for HDF5's broader out-of-space
write, flush, and close problem. \\

FFM / JDK & Foreign Function and Memory API / Java Development Kit &
The JDK supplies the Java compiler and runtime tools; its FFM API lets Java
code call native libraries and access foreign memory. HDF5 can optionally
generate Java FFM bindings. \\

GHSA & GitHub Security Advisory &
A security-advisory record and identifier published through GitHub. Like a CVE,
it must be checked for exact scope, evidence, and applicability. \\

HDF / HDF5 & Hierarchical Data Format / Hierarchical Data Format version 5 &
In this report, HDF5 refers to the version-5 format and data model and, only
where specifically stated, the core C library, official CLI tools, extension
classes, and build or release configuration. HDF4 and older HDF technology are
different formats and software. \\

H5II / H5PW & HDF5 Independent Implementation / HDF5 Python Wrapper &
Names of two SSP SIG mechanism models used in the independent-implementation
and binding chapters. Their numbered families classify mechanisms rather than
assigning this report's risk tiers. \\

HDFS & Hadoop Distributed File System &
A distributed storage system that can be reached through an optional HDF5 VFD
and therefore adds build, storage, and deployment dependencies. \\

HPC & High-Performance Computing &
Computing environments such as clusters and supercomputers, often involving
parallel I/O, MPI, shared storage, and tightly matched software builds. \\

HSDS & Highly Scalable Data Service &
An HDF5-oriented service shown for ecosystem context but expressly outside the
assessment scope. \\

HTTP & Hypertext Transfer Protocol &
A network protocol used to retrieve resources. HTTP requests, range requests,
proxies, and access logs can expose data locations and access patterns. \\

ID & Identifier &
A name or code used to distinguish a record, object, process resource, or other
entity. Report identifiers and HDF5 object identifiers belong to different
namespaces and must be interpreted in context. \\

I/O & Input / Output &
Reading or writing data between a process and files, memory, devices, or
network services. \\

JNI / JVM & Java Native Interface / Java Virtual Machine &
JNI connects Java code to native libraries; the JVM executes Java bytecode.
Their boundary matters when Java-facing software reaches native HDF5 code. \\

LSan & LeakSanitizer &
A runtime testing tool that detects allocated memory that becomes unreachable
without being released. \\

LRU & Least Recently Used &
A cache-management policy that preferentially evicts entries not used recently.
The report mentions LRU rank in one metadata-cache invariant case. \\

MPI / MPI-I/O & Message Passing Interface / MPI input and output &
Standards used for parallel programs and coordinated parallel file access.
Parallel HDF5 builds depend on compatible MPI behavior and binaries. \\

N/A & Not applicable &
A report status for an item that is not a risk-bearing finding, such as a
recommendation or program theme. It does not mean that a risk was evaluated as
zero. \\

\code{NDEBUG} / \code{-DNDEBUG} & Assertions disabled &
The C/C++ convention and compiler definition that remove ordinary
\code{assert(...)} checks. This commonly occurs in release builds. \\

NE & Not evaluated &
A rating status used when impact or exposure lacks enough evidence. NE never
means Low risk. \\

OSS & Open-Source Software &
Software whose source is available under an open-source license. Openness
supports inspection and reuse but is not, by itself, a security guarantee. \\

OSS-Fuzz & Continuous fuzzing service for open-source software &
A service that repeatedly runs project fuzz targets at scale. Its results
cover the submitted targets and configurations rather than every possible HDF5
path. \\

POSIX & Portable Operating System Interface &
A family of operating-system interface standards. POSIX calls and semantics
underlie some HDF5 file drivers and candidate storage-reservation controls. \\

PR & Pull Request &
A proposed source change submitted for review and automated checks before
merging. A PR test gate provides fast feedback but not long-running assurance. \\

RCE & Remote Code Execution &
Execution of attacker-chosen code across a remote trust boundary. The broader
phrase \emph{arbitrary code execution} does not necessarily imply a remote
attacker. \\

REST & Representational State Transfer &
A common style of network API that some connectors or external services may
use. \\

ROS3 & Read-Only S3 &
The HDF5 Virtual File Driver for read-only access to objects through an S3 or
S3-compatible interface. Its network, credential, and service behavior remain
part of the deployment boundary. \\

SBOM & Software Bill of Materials &
A machine-readable inventory of software components and versions in a build or
package. It supports applicability analysis but does not prove that a
vulnerability is reachable or fixed. \\

S3 & Amazon Simple Storage Service, or an S3-compatible interface &
An object-storage API used by cloud services and by HDF5's ROS3 VFD for
read-only access. The governing provider need not be Amazon when an
S3-compatible service is used. \\

SDLC & Software Development Life Cycle &
The processes used to design, implement, test, release, operate, and maintain
software. \\

SIG & Special Interest Group &
A community group focused on a defined subject. The SSP SIG is the HDF5 Safety,
Security \& Privacy Special Interest Group. \\

SIGFPE & Floating-point exception signal &
An operating-system signal commonly raised by invalid integer or floating-point
arithmetic, including the division-by-zero case discussed in the report. \\

SLA & Service-Level Agreement &
A documented operating commitment, such as a response target. The proposed
roadmap calls for SLA child records to make clocks and evidence accountable;
those records do not create a risk rating or show that the proposal was
adopted. \\

SLSA & Supply-chain Levels for Software Artifacts &
A framework for improving and describing the integrity and provenance of
software builds and releases. \\

SSP & Safety, Security \& Privacy &
The three assessment dimensions and the name of the HDF5 SSP SIG. Their
meanings in this report are separately defined below. \\

SWMR & Single Writer / Multiple Reader &
An HDF5 access mode that permits one writer and multiple readers to access a
file concurrently under SWMR's required open, flush, and refresh protocol and
filesystem semantics. It is not a general crash-recovery or multi-writer
guarantee. \\

\code{T0} / TBD & Clock start / To be determined &
\code{T0} is the recorded start date for a finding-and-scope or work-package
clock. Depending on its stated trigger, it is affected-scope or equivalent-
exposure identification, or roadmap adoption; the record must say which. TBD
marks a field for which a required decision has not yet been made. \\

UBSan & Undefined Behavior Sanitizer &
A runtime testing tool that detects selected C/C++ operations whose behavior
the language does not define, such as some invalid arithmetic. \\

UAF & Use After Free &
A memory-lifetime error in which code accesses an object after releasing its
storage. The full term is defined in the concepts below. \\

URL & Uniform Resource Locator &
An address for a network resource. URLs in metadata, logs, or connector
configuration can themselves reveal sensitive context. \\

VDS & Virtual Dataset &
An HDF5 dataset whose logical view maps to regions of source datasets, possibly
in other files. Reading it can therefore cause external-file access. \\

VFD & Virtual File Driver &
The HDF5 layer that maps logical file reads and writes to physical storage,
such as a local file, memory, multiple files, MPI-I/O, or object storage. \\

VLEN & Variable Length &
An HDF5 datatype whose elements have variable size. VLEN data require dynamic
allocation and careful lifetime and resource handling. \\

VOL & Virtual Object Layer &
The HDF5 layer through which a connector implements object operations. A VOL
connector can present remote, generated, or non-HDF5 data through the HDF5 API. \\

\end{longtable}
\endgroup

\subsection{Report identifiers and notation}

Identifiers are labels, not scores. A recommendation number does not imply
priority, and a roadmap phase does not change a finding's risk. The letter
\code{I} is context-sensitive: \code{I1}--\code{I4} mean impact levels in a
risk record, while \code{I1}--\code{I9} in the independent-implementation
chapter name mechanism families. The surrounding table or model determines
which meaning applies.

\begingroup
\small
\setlength{\tabcolsep}{3pt}
\renewcommand{\arraystretch}{1.12}
\begin{longtable}{@{}>{\raggedright\arraybackslash}p{0.25\linewidth}p{0.69\linewidth}@{}}
\hline
\textbf{Pattern or label} & \textbf{Meaning} \\
\hline
\endfirsthead
\hline
\textbf{Pattern or label} & \textbf{Meaning} \\
\hline
\endhead
\hline
\endfoot

\code{CORE}, \code{APP}, \code{EXT}, \code{FMT}, \code{IND},
\code{LONG}, \code{SUP} &
Report component prefixes for the core library or tools, downstream
applications, extensions, format governance, independent implementations,
long-term accessibility, and the software supply chain. \\

\code{SEC}, \code{SAFE}/\code{SAF}, \code{PRIV}, \code{TOOL} &
Report dimension or surface labels for security, safety, privacy, and an
official tool. \\

\code{<component>-<dimension>-nnn} &
A report finding or evidence-record identifier, for example
\code{CORE-SEC-001}. The register states whether the record is formally rated
or remains in the evidence backlog. \\

\code{<component>-Snn} &
A recommendation identifier, for example \code{APP-S4}. Variants such as
\code{FILTER-S1}, \code{VFD-S4}, and \code{CORE-PRIV-S2} identify a more
specific track. The number is an identifier only. \\

\code{RM-nX} &
A phased-roadmap work package. Its phase and sequence describe delivery
planning, not severity or inherent risk. \\

\code{I1}--\code{I4} and \code{E1}--\code{E4} &
The report's ordinal impact and exposure levels. They are combined through the
matrix in Table~\ref{tab:inherent-risk-matrix}; they are not multiplied as
numbers. \\

\code{P1}--\code{P7} &
Privacy-exposure mechanism families: semantic metadata; structural or on-disk
leakage; raw-data patterns; unintended inclusion during processing;
aggregation, correlation, or inference; persistence beyond the intended
lifetime; and misplaced trust in privacy by format. They are categories, not
risk ratings. \\

\code{W1}--\code{W10} &
Wrapper or binding mechanism families from the SSP SIG H5PW model. They cover
native compromise, code in data, dynamic loading, lifecycle mismatch, resource
amplification, ABI or packaging drift, artifact exposure, trust-boundary
confusion, misleading safety signals, and external-reference traversal. \\

\code{I1}--\code{I9} in the H5II model &
Independent-implementation mechanism families covering parser/specification
divergence through trust-boundary confusion. These are unrelated to the
report's four impact levels despite reusing the letter \code{I}. \\

\code{ATK-nnn}, \code{HAZ-nnn}, \code{EXP-nnn} &
SSP SIG threat, safety-hazard, and privacy-exposure registry identifiers. They
are separate program artifacts, not findings in this report. \\

\code{EXT}, \code{OPS}, \code{TCD}, \code{SCD} in SSP tag fields &
SSP classification tags for extensions, operations, toolchains, and
distribution, respectively. They classify a mechanism and carry no report-wide
risk rating. \\

\code{HDF5-SHINES-}\allowbreak\code{FAMILY-nn} &
An SSP SIG policy-control identifier. A crosswalk to a report recommendation
does not establish adoption or control satisfaction. \\

\code{APP}, \code{ASSURE}, \code{CORE}, \code{EXT}, \code{GOV},
\code{INTEROP}, \code{PRIV}, \code{REL} in owner fields &
Accountability role labels for application, assurance, core engineering,
extension, governance, interoperability, privacy, and release work. They name
roles rather than currently assigned people or organizations. \\

\end{longtable}
\endgroup

\subsection{Quick distinctions}

Two groups of terms recur in conclusions and recommendations but answer
different questions. These tables are reading aids; the full definitions and
the controlling rules remain in the cited report sections.

\subsubsection{Risk and response terms}

\begingroup
\small
\setlength{\tabcolsep}{3pt}
\renewcommand{\arraystretch}{1.12}
\begin{longtable}{@{}>{\raggedright\arraybackslash}p{0.20\linewidth}>{\raggedright\arraybackslash}p{0.31\linewidth}>{\raggedright\arraybackslash}p{0.43\linewidth}@{}}
\hline
\textbf{Term} & \textbf{Question it answers} & \textbf{Use in this report} \\
\hline
\endfirsthead
\hline
\textbf{Term} & \textbf{Question it answers} & \textbf{Use in this report} \\
\hline
\endhead
\hline
\endfoot

Impact & How serious is the greatest credible consequence? &
Ordinal level \code{I1}--\code{I4}, supported by evidence for the bounded
scenario. \\

Exposure & How readily and commonly can the scenario arise? &
Ordinal level \code{E1}--\code{E4}, based on reachability, preconditions,
deployment prevalence, and recurrence; not a numerical probability. \\

Inherent risk and risk tier & What does the report rate before optional or
proposed controls? & The impact/exposure matrix yields Low, Moderate, High, or
Critical for a bounded scenario. The two ordinal levels are not multiplied. \\

Rating status & Is the evidence sufficient to issue that tier? &
Rated and Provisional records carry a tier; NE means a required input is not
evaluated sufficiently, and N/A marks a non-risk-bearing item. \\

Current residual risk & What remains after controls already verified in the
stated deployment? & It can differ from inherent risk only when evidence shows
that existing controls apply and work. \\

Target residual risk & What might remain if proposed controls are completed? &
A projection, not achieved risk reduction; acceptance evidence must pass before
it can become current residual risk. \\

Response urgency & How quickly does the report recommend treatment? &
A report-assigned route and target horizon based on remaining risk, time to
harm, recurrence or exploitation, and risk tolerance; it is not an adopted
stakeholder deadline. \\

SSP SIG severity and likelihood & How does a separate SSP registry score a
mechanism? & Separate 1--5 inputs that the SSP SIG multiplies. Its mechanism
scores cannot be converted into or merged with this report's scenario ratings. \\

\end{longtable}
\endgroup

\subsubsection{Format interpretation terms}

\begingroup
\small
\setlength{\tabcolsep}{3pt}
\renewcommand{\arraystretch}{1.12}
\begin{longtable}{@{}>{\raggedright\arraybackslash}p{0.24\linewidth}>{\raggedright\arraybackslash}p{0.70\linewidth}@{}}
\hline
\textbf{Term} & \textbf{What it establishes} \\
\hline
\endfirsthead
\hline
\textbf{Term} & \textbf{What it establishes} \\
\hline
\endhead
\hline
\endfoot

Conformance & The bytes or implementation satisfy the normative requirements
of the governing specification. \\

Reader acceptance & One particular reader, version, build, and profile opens
or processes a file. It does not alone establish conformance. \\

Compatibility & A stated producer, artifact, reader, and configuration work
together in practice. It is an observed relationship, not a universal property. \\

Backward / forward compatibility & Respectively, newer software handles older
artifacts, or older software handles newer artifacts within a declared feature
set. Neither licenses malformed metadata. \\

Interoperability & Different tools, versions, or implementations exchange and
interpret data as intended across a stated scope. \\

Deployment safety & Processing the artifact under a deployment's authority,
limits, and policy does not produce the adverse behavior in question. A file
can conform and still be unsafe for a particular deployment. \\

\end{longtable}
\endgroup

\subsection{Terms and concepts}

PDF search is usually the quickest lookup method. The entries are also grouped
in four navigable ranges: \hyperref[glossary:ac]{A--C},
\hyperref[glossary:dh]{D--H}, \hyperref[glossary:ip]{I--P}, and
\hyperref[glossary:qz]{Q--Z}. A slash in a headword groups terms that are most
usefully distinguished together; common secondary terms are repeated or
cross-referenced where alphabetical lookup would otherwise be difficult.

\subsubsection{A--C}
\label{glossary:ac}

\begin{description}[leftmargin=1.5em,style=nextline]

\item[Accountability roles.]
The accountable owner answers for an outcome; the implementer performs the
work; the verifier independently reproduces acceptance evidence; the approver
authorizes rollout; and the risk acceptor decides whether verified remaining
risk may be retained. The roadmap separates these roles for High and Critical
scope.

\item[Acceptance criteria / acceptance evidence.]
The observable condition, test result, document, or operating record required
to show that a recommendation was implemented and works over its stated scope.
Completing an activity is not acceptance evidence unless its required outcome
is also demonstrated.

\item[Access control.]
Rules and mechanisms that decide which identities may read, change, create, or
delete a resource. HDF5 applications normally rely on the operating system,
storage service, or surrounding application for identity and access control.

\item[Adversary / attacker / malicious actor.]
A person or system intentionally trying to cause an unwanted result. Safety
failures and privacy harms in this report can also occur without an adversary.

\item[Adverse-event scenario.]
A bounded statement connecting an initiating actor or non-adversarial event,
an assessment object, exposure and preconditions, a condition, and a credible
harm to named assets, people, or operations. This---not a technology or defect
name alone---is what the report rates.

\item[Affected scope / affected population.]
The specific versions, builds, configurations, deployments, files, users, or
workflows to which a scenario applies. A project name or CVE count is not a
substitute for establishing this population.

\item[Aggregation, correlation, and inference.]
Combining or analyzing data and metadata so that the result reveals information
not evident from any one item. This can create privacy harm even when every
individual input was disclosed legitimately.

\item[Allocation / resource budget.]
Allocation reserves memory, storage, or another finite resource. A resource
budget is an enforced upper bound on aggregate or per-object use, such as
logical bytes, objects, traversal depth, CPU time, or decompression expansion.

\item[Allowlist.]
An explicit list of permitted items, such as approved plugin identities or
external directories. Items not on the list are denied. An allowlist is useful
only if its storage and update path are themselves protected.

\item[Ambient authority.]
Permissions already held by a process---for example filesystem, network,
credential, or plugin-loading access---that are available without a fresh,
resource-specific authorization decision. An untrusted file can be dangerous
when its metadata directs a process to exercise this authority.

\item[Application boundary.]
The point where an application decides how an HDF5 capability will be used and
with what permissions. A defect at this boundary is not automatically a defect
in the HDF5 parser merely because HDF5 carried the instruction.

\item[Archive / archival profile.]
An archive preserves information for later use. An archival profile states the
formats, features, dependencies, and validation rules accepted for that purpose;
preserving primary file bytes alone may be insufficient.

\item[Artifact.]
A durable output used as evidence or input. Depending on context this can mean a
source or release package, a test fixture, an SBOM, a signed manifest, a log or
cache created by a workflow, or a report and decision record.

\item[Assertion / always-active check.]
An assertion is a developer check written with facilities such as
\code{assert(...)} and often removed from release builds. An always-active
runtime check remains present in supported production configurations. A
condition derived from untrusted file bytes must not depend on an assertion
alone.

\item[Assessment horizon.]
The period over which a scenario is evaluated. Time is material for risks such
as a filter becoming unavailable over a decade; ratings with materially
different horizons should not be directly ordered without qualification.

\item[Assessment object.]
The exact technology, version, file population, component, or deployment being
examined. A formal rating attaches to this bounded object and scenario, not to
``HDF5'' in the abstract.

\item[Assessment surface / scope level.]
A surface is a component, input, interface, or boundary considered while
looking for risk. A \emph{primary assessment object} is directly assessed; a
\emph{boundary case study} examines only a defined interaction with a primary
object; and a \emph{contextual reference} illustrates the landscape without
creating an audit conclusion about that product. Appearing in a diagram,
example, or citation does not mean the entire product was audited.

\item[Assurance.]
Justified confidence, based on stated evidence, that a property holds over a
defined scope. Tests, review, reproducible builds, and independent verification
can increase assurance, but assurance is not an absolute guarantee and does not
by itself mean certification or an assurance opinion.

\item[Asset.]
Something of value that the assessment seeks to protect, such as scientific
data, metadata, credentials, software integrity, service availability,
research validity, privacy, or long-lived records.

\item[Attack mechanism.]
A capability or sequence an attacker can use to produce harm. A documented and
correctly implemented feature can be an attack mechanism when a surrounding
application grants it inappropriate authority.

\item[Attack surface.]
The set of inputs, interfaces, features, tools, and extension points through
which an adverse condition might be reached. More surface does not by itself
establish a vulnerability or a risk tier.

\item[Attestation / certification / assurance opinion.]
Formal statements made under defined criteria and evidence requirements. This
technical audit is not any of these and does not certify an HDF5 release,
application, file, or deployment.

\item[Audit.]
A documented technical examination in this report. It does not mean a
compliance audit, certification, attestation, penetration test, or exhaustive
line-by-line source review.

\item[Authentication / authorization.]
Authentication establishes which identity is making a request; authorization
decides what that identity may do. Successfully authenticating a connector or
plugin does not authorize every file-selected action it can perform.

\item[Authority.]
The practical ability to act on a resource, such as reading a local file,
opening a network connection, allocating memory, loading code, or using a
credential. The report assigns a control failure to the boundary that supplies
and governs the relevant authority.

\item[Availability.]
The property that data and services remain accessible and usable when needed.
Confidentiality, integrity, and availability are the conventional three core
security objectives.

\item[Backport.]
Applying a correction developed for a newer code line to an older maintained
release. A backport is not assumed to contain the fix until the exact release
artifact and path are verified.

\item[Backward / forward compatibility.]
Backward compatibility lets newer software use older artifacts; forward
compatibility lets older software handle artifacts from newer producers within
a defined feature set. Neither requires accepting malformed metadata, and
historical reader acceptance does not prove format conformance.

\item[Baseline.]
The exact version, commit, build, configuration, or deployment state against
which an observation is made. This historical ecosystem assessment has no
single baseline that represents every conclusion.

\item[Binary / source / release artifact.]
Source is human-readable program material; a binary is compiled executable
code; a release artifact is a published package containing either or both.
Integrity and provenance must identify which exact artifact was assessed.

\item[Binding / wrapper.]
Software that exposes HDF5 to another language or programming model, such as
Python, Java, or .NET. It translates types, memory ownership, errors, and
lifecycle rules, but it does not automatically make every HDF5 capability safe
for an application.

\item[Bounded decode.]
Reading a defined number of bytes through a cursor or equivalent boundary so a
decoder cannot silently run past the available record. The report's preferred
sequence is decode, validate, construct, and only then separately authorize
activation.

\item[Buffer overflow / out-of-bounds read or write / buffer over-read.]
An overflow or out-of-bounds write stores data beyond the memory reserved for a
buffer; an out-of-bounds read or over-read reads beyond it. Either can crash a
process, disclose memory, corrupt state, or in some conditions enable code
execution.

\item[Build configuration / build profile.]
The selected compiler, options, dependencies, features, and debug or release
settings used to produce a binary. Different HDF5 builds of the same nominal
version may have materially different behavior and dependencies.

\item[Cache / temporary artifact.]
A cache retains data to avoid recomputation or retrieval; a temporary artifact
is an intermediate file, buffer, snapshot, or log. Either may persist longer
than intended and expose data outside the primary HDF5 file.

\item[Callback.]
A function or handler supplied for another component to invoke later, often
when an event or extension point is reached. A callback executes code in the
calling process, so its origin, lifetime, inputs, and authority are part of the
security boundary.

\item[Canonical / minimal / nonminimal encoding.]
A canonical encoding is the one representation selected when a rule permits
several byte representations of the same value. A minimal encoding uses the
shortest sufficient representation; a nonminimal encoding uses more space. A
nonminimal encoding can still conform unless the specification expressly
requires minimal or canonical form.

\item[Capability.]
An ability a component exposes or holds, such as opening an external path,
contacting a service, allocating memory, or loading executable code. A
capability may be intentional and correctly implemented yet still require
explicit authorization and limits at a trust boundary; it is not synonymous
with a vulnerability.

\item[Census.]
A systematic inventory of installed versions, linkage, features, files, or
deployments. The roadmap census supplies affected-scope evidence; discovering
scope is not itself treatment.

\item[Checked arithmetic.]
Addition, multiplication, range calculation, or conversion that detects
overflow, underflow, division by zero, truncation, or an invalid bound before
using the result for allocation, copying, or traversal.

\item[Checkpoint.]
A known recovery point recorded during a long-running write or computation.
Checkpoint discipline limits how much work or data must be reconstructed after
a failure but is not the same as a file-format transaction guarantee.

\item[Checksum / cryptographic hash / digital signature.]
A checksum helps detect accidental change. A cryptographic hash produces a
substitution-resistant digest, but detects malicious replacement only when the
expected digest comes through a trusted or authenticated channel. A verified
digital signature can provide integrity and origin authentication when the
verification key is reliably associated with the signer. None of these
mechanisms proves that the checked code is safe.

\item[Chunk / chunked storage / chunk index.]
A chunk is a separately stored block of a dataset. Chunking supports partial
I/O, compression, and extendible datasets. A chunk index is the on-disk
structure that maps logical chunk coordinates to stored blocks.

\item[Ciphertext / plaintext.]
Plaintext is information in a directly readable form; ciphertext is the output
of encryption. Data become plaintext again during authorized processing, where
applications, plugins, caches, and logs can expose them.

\item[Closure: defect, class, recommendation, and risk.]
Defect closure shows that one affected case was treated. Class closure shows
that the recurring family and sibling paths are structurally covered.
Recommendation closure requires its acceptance evidence, while risk closure
also requires an accountable decision about what remains. These claims are not
interchangeable.

\item[Codec.]
Software that encodes and decodes a data representation, such as a compression
filter. Long-term readability depends on retaining either a working codec or
enough specification and test evidence to reproduce it.

\item[Command-line tool.]
A standalone program invoked from a shell. HDF5 inspection and conversion tools
are security and privacy surfaces because they parse files, render output, load
extensions, and can create persistent copies.

\item[Compatibility.]
The practical ability of particular producers, files, readers, and
configurations to work together. Compatibility is an observed relationship and
is distinct from both conformance to the specification and deployment safety.

\item[Compensating control.]
A control used when the preferred or structural control is unavailable, such
as isolating an affected parser until an upgrade can be deployed. Its coverage
and effectiveness must be tested rather than assumed.

\item[Condition / failure mode / consequence.]
The condition is the state that permits an event, the failure mode is how the
component departs from intended behavior, and the consequence is the resulting
effect. Separating them prevents a visible crash from being mistaken for the
earliest root cause.

\item[Confidentiality.]
The property that information is not disclosed to unauthorized recipients.
Privacy is broader: inference, retention, correlation, or unwanted secondary
use can cause privacy harm even without a confidentiality breach.

\item[Conformance.]
Compliance with the normative requirements of a specification. A file can be
accepted by one reader yet be non-conforming, or be conforming and rejected by
a reader that imposes an extra condition.

\item[Confused deputy.]
A system with legitimate authority that is induced to use that authority on
behalf of a less-privileged requester. In this report, a file-processing service
that follows an attacker-selected local path with the service's permissions is
the central example.

\item[Connector / connector stack.]
An extension implementing an HDF5 abstraction, especially a VOL connector.
Multiple connectors may be layered as a stack, so data origin, caching,
network use, and side effects may span more than one component.

\item[Containment / structural prevention.]
Containment promptly limits a known exposure, for example by disabling a
feature, isolating a process, or monitoring capacity. Structural prevention
changes the design, implementation, or process so the recurring defect class
is less likely to return. One does not substitute for the other's deadline.

\item[Control.]
A technical, operational, or governance measure intended to prevent, detect,
limit, or recover from an unwanted event. A proposed control receives credit
in residual risk only after relevant acceptance evidence exists.

\item[Controlled rejection.]
Stopping on invalid or disallowed input with a bounded, documented diagnostic,
without a crash, uncontrolled allocation, unintended side effect, or silent
repair. Controlled rejection is often the safe result for unsupported features.

\item[Core VFD.]
The specifically named in-memory Virtual File Driver, optionally with a backing
file written at close. In that phrase, \emph{Core} does not mean the HDF5 core
library.

\item[Corpus / fixture / test vector.]
A corpus is a maintained collection of test inputs. A fixture is one input with
a stated expected result; a test vector pairs known input and output. Positive
fixtures should be accepted, negative fixtures should be rejected, and
contested fixtures preserve competing interpretations until adjudication.

\item[Coverage.]
Evidence about which declared code, structures, invariants, configurations, or
cases a test or control exercised. A percentage is meaningful only with a
versioned denominator and does not, by itself, prove defect or class closure.

\item[Crash consistency.]
The guarantees about durable file state when a writer crashes, is killed, or
loses power during an update. A flush call, SWMR mode, or successful earlier
write does not by itself define those guarantees.

\item[Credential / secret.]
Information that grants authority, such as a token, password, private key, or
cloud credential. Logs, environment variables, crash dumps, connectors, and
plugins can expose secrets even when HDF5 payload data are protected.

\item[Current working directory / search path.]
The current working directory is the process's default relative-path base. A
search path is the ordered set of locations examined for an external file or
loadable module. User-writable or implicit locations can let untrusted parties
control what is opened or executed.

\end{description}

\subsubsection{D--H}
\label{glossary:dh}

\begin{description}[leftmargin=1.5em,style=nextline]

\item[Data integrity.]
The property that data remain correct, complete, and unaltered except through
authorized changes. In this report it includes avoiding silent wrong data,
corruption, and scientifically invalid interpretation.

\item[Data lifecycle.]
The stages through which data pass, including creation, processing, storage,
sharing, transfer, reuse, archival, and deletion. Security and privacy controls
must account for artifacts created at every relevant stage.

\item[Data minimization / tokenization.]
Minimization avoids or removes information that is not needed. Tokenization
replaces a sensitive value with a less revealing reference. A field is not
actually minimized if it disappears from the API view but remains recoverable
from raw bytes, free space, logs, or copies.

\item[Data remanence / secure deletion.]
Data remanence is information that remains recoverable after ordinary deletion
or replacement, including in free blocks, histories, replicas, snapshots, and
backups. Secure deletion is a storage-specific process intended to make the
defined information unrecoverable across the stated copies and layers.

\item[Dataset.]
An HDF5 object containing an array of data values together with a datatype and
dataspace, plus optional attributes and storage properties.

\item[Dataspace.]
A dataspace defines a dataset's or attribute's extent: scalar, null, or simple,
with rank and current or maximum dimensions. A dataspace handle can also carry
a selection used for dataset I/O; that selection is not normally part of the
stored dataset or attribute definition. A declared logical extent can be much
larger than the physically stored payload.

\item[Datatype / committed datatype.]
The description of how values are represented and interpreted. A committed
datatype is stored as a named HDF5 object that other objects can reference;
compound, variable-length, enumeration, and nested types add parsing complexity.

\item[Decode / encode.]
Decode turns stored bytes into an in-memory representation; encode performs the
reverse. The report requires bounded decode followed by record-, object-, and
graph-level validation before construction; invalid state must not reach later
encode or flush paths.

\item[Default deny.]
A policy under which a capability is unavailable unless it is explicitly
permitted. For untrusted HDF5 input, external resources and executable plugins
are examples of capabilities suited to a default-deny policy.

\item[Defect / weakness / vulnerability.]
A defect is a specific implementation or specification error. A weakness is a
broader condition or recurring class that can contribute to harm. A
vulnerability is a weakness that can be reached and used under stated
conditions to violate a security property. The words are related but are not
interchangeable, and none alone assigns a report risk tier.

\item[Defense in depth.]
Using independent layers of protection so one failure does not immediately
produce harm, such as combining input validation, resource limits, process
isolation, and least privilege.

\item[Denial of service / resource exhaustion / amplification.]
Denial of service makes a system unavailable. Resource exhaustion consumes
finite memory, CPU, storage, output, or handles. Amplification occurs when a
small input induces disproportionately large work or allocation.

\item[Dependency / upstream / downstream / transitive dependency.]
A dependency is software or infrastructure required by another component.
Upstream dependencies are used to build or run HDF5; downstream projects use or
package HDF5. A transitive dependency is reached indirectly through another
dependency, which can make the actual HDF5 build invisible to an end user.

\item[Deployment.]
A concrete installed system: its versions, build options, enabled features,
configuration, permissions, users, data sources, storage, and surrounding
controls. Risk conclusions often depend more narrowly on a deployment than on
a product name.

\item[Deserialization.]
Reconstructing application objects or behavior from stored data. HDF5 decoding
can be safe while a higher-level model loader unsafely reconstructs executable
objects from HDF5 content.

\item[Differential testing / differential conformance testing.]
Giving the same inputs to multiple builds or independent implementations and
comparing outcomes. A difference exposes a question; it does not by itself
prove which implementation or specification is wrong.

\item[Digital signature.]
See \emph{Checksum / cryptographic hash / digital signature}. A signature is
useful only when verification succeeds and the verification key is reliably
associated with the claimed signer.

\item[Disposition / ruling.]
An authoritative resolution of a reported issue or specification question,
including which interpretation governs and which artifacts or implementations
must change. An issue label, implementation behavior, or proposed patch is not
a ruling unless the responsible authority records it as such.

\item[Documented observation / reasoned inference.]
A documented observation is supported directly by a cited artifact, with
source-reported evidence distinguished from independent reproduction. A
reasoned inference draws a conclusion from evidence plus stated assumptions
and uncertainty. Neither is a recommendation or a formal rating by itself.

\item[Durability.]
The property that data acknowledged as written survive the failures covered by
the stated guarantee. Durability differs from an in-memory value being visible
to a reader and must be evaluated against the storage stack and recovery model.

\item[Dynamic loading / plugin loading.]
Loading executable code while a process is running, often from a shared
library found on a search path. The loaded code executes with the process's
authority, so its identity, origin, path, configuration, and privileges matter.

\item[Effort / implementation dependency / sequence.]
Effort estimates the scale of delivery work; a dependency is a concrete
prerequisite; sequence records an implementation order and its rationale.
These planning fields may shape delivery but never lower or otherwise change
inherent risk or response urgency.

\item[Embargo.]
A temporary restriction on vulnerability information while affected parties
coordinate analysis, fixes, and disclosure. Embargoed local repair can close a
known defect without proving that the recurring class has been eliminated.

\item[Encryption / key management.]
Encryption uses a key to make protected information unreadable at a stated
boundary. Key management covers generation, custody, access, backup, rotation,
recovery, and destruction. Encryption that leaves metadata or keys exposed, or
whose keys cannot be recovered, does not satisfy the report's broader privacy
and durability needs.

\item[Error path / cleanup / unwinding.]
The code followed after an operation fails, including releasing partially
created objects and restoring state. Malformed input often exposes lifetime
defects in these paths even when the original decode detects an error.

\item[Evidence action / evidence backlog.]
A bounded investigation needed because a scenario lacks enough information to
rate or route. A backlog record is not a formal risk tier and must never be
read as Low risk.

\item[Evidence confidence.]
The report's judgment about how directly and completely the evidence supports
a scenario and its rating assumptions. High, Moderate, or Low confidence is
reported separately from impact, exposure, and risk tier.

\item[Evidence cutoff / subsequent-event update.]
The cutoff is the latest date on which substantive evidence is eligible for
this revision. A later retrieval can document an older artifact, but a new
measurement, fix, or decision after the cutoff is a subsequent event unless
the report expressly incorporates it.

\item[Exception / variance.]
An approved, scoped departure from a control, gate, route, or deadline. The
roadmap requires it to name an accountable authority, rationale, interim
control, expiration, and remaining-risk decision; an undocumented deviation is
not an exception record.

\item[Exploitability / reachability.]
Reachability asks whether a real input can enter the relevant code path in a
specified build and deployment. Exploitability asks whether the resulting
behavior can be used to cause a security consequence. A source pattern or
crash alone may establish neither for supported releases.

\item[Exposure (risk axis).]
The E1--E4 axis expressing how readily and commonly a stated scenario arises,
considering reachability, preconditions, deployment prevalence, and recurrence.
It is ordinal and is not a numerical probability. \emph{Privacy exposure} is a
different use of the word: information becoming observable or usable in a way
that may cause privacy harm.

\item[Extension.]
A component that adds behavior outside the baseline library path, including a
filter, VOL connector, or VFD. An extension may be built in or loaded as a
plugin; the words \emph{extension} and \emph{dynamically loaded plugin} are
therefore not synonyms.

\item[External link / external raw storage / external reference.]
An external link names an object in another HDF5 file. External raw storage
places a dataset's raw bytes in non-HDF5 files. VDS mappings can also name
external source files. This report uses \emph{external reference} or
\emph{external resource} as an umbrella for file-controlled access beyond the
initial file; it is not limited to the HDF5 object-reference datatype.

\item[Fail closed / fail open.]
A fail-closed system denies or safely stops when validation, policy, or a
dependency cannot be established. A fail-open system continues, guesses, or
silently ignores the condition. Unsupported or unknown untrusted-file features
should normally fail closed.

\item[Fault injection.]
Deliberately causing failures---such as out-of-space errors, failed writes, or
process termination---at controlled points to measure durable state, error
reporting, and recovery behavior.

\item[File-derived value / file-controlled metadata.]
A size, count, address, path, shape, filter setting, link, flag, or other value
obtained from the input file. Until validated and authorized, it must be treated
as controlled by the file's originator rather than trusted program state.

\item[File disclosure / data exfiltration.]
File disclosure is unauthorized reading or return of file contents.
Exfiltration is their unauthorized transfer out of the protected boundary.
The demonstrated external-storage incident in this report is rated for the
disclosure actually evidenced, not for an unreported code-execution path.

\item[File format / data model / library.]
The file format specifies the on-disk representation; the data model describes
logical objects and relationships; the library is software that reads, writes,
and manipulates them. A defect or capability in one is not automatically a
defect in the others.

\item[File locking.]
Coordination intended to prevent incompatible concurrent access to a file.
Its availability and semantics depend on the HDF5 build, configuration,
filesystem, and access pattern; it is not a substitute for crash recovery.

\item[Fill value.]
The value HDF5 returns for dataset elements that have not been explicitly
written, subject to the dataset's creation and allocation properties. A fill
value is part of dataset behavior and should not be confused with bytes that
happen to remain in unused storage.

\item[Filter / filter pipeline / filter plugin.]
A filter transforms dataset chunks during writing and reverses or interprets
the transformation during reading. A pipeline is an ordered sequence of
filters. A plugin supplies filter code dynamically; losing that dependency can
make data unreadable.

\item[Finding / formal finding.]
A documented adverse-event scenario supported strongly enough to enter the
authoritative findings register. Its identifier, scope, evidence, rating, and
response route are distinct from the recommendations that may address it.

\item[Flush.]
An operation that pushes buffered HDF5 state toward a lower storage layer.
What becomes durable after a flush depends on the documented library, driver,
operating-system, and storage guarantees; the word alone does not promise
crash consistency.

\item[Format bounds / library-version bounds.]
HDF5 properties that constrain which on-disk format versions and features a
writer may emit. They are not the same as the installed library's version and
do not alone prove reader compatibility.

\item[Fuzzing / mutation.]
Automated testing that generates or alters many inputs to find crashes,
resource problems, or invariant violations. Whole-file fuzzing mutates bytes
broadly; structure-aware or invariant-guided fuzzing targets meaningful HDF5
records and boundary values.

\item[Graph / traversal.]
HDF5 files contain linked structures that can be viewed as a graph. Traversal
follows those relationships. Cycles, excessive depth, repeated work, or
external targets require explicit progress rules and budgets.

\item[Group.]
An HDF5 container object whose links organize datasets, committed datatypes,
and other groups into a hierarchy, much like directories organize filesystem
names.

\item[Guard / dominating check.]
A guard is a check that prevents unsafe execution when its condition fails. A
check dominates an operation when every path to that operation must pass
through the check. Cataloging a check matters because a nearby assertion may
not protect all release-build paths to the unsafe use.

\item[Hard link / soft link.]
An HDF5 hard link is a group entry that names an object in the same file; more
than one hard link can name the same object. A soft link stores a path to be
resolved later and can be dangling. An external link additionally names another
file and crosses an external-resource boundary.

\item[Hazard.]
A condition with the potential to cause a non-adversarial safety loss. The SSP
SIG hazard registry and this report use different rating methods, so a hazard
identifier or its registry score is not a report finding or tier.

\item[HDF Group, The.]
The nonprofit organization that develops and supports the reference HDF5
software and stewards the format. It is distinct from the broader ecosystem of
downstream and independent projects.

\item[HDF5 interface and subsystem prefixes.]
Prefixes such as \code{H5C}, \code{H5O}, \code{H5T}, and \code{H5Z} identify
metadata-cache, object, datatype, and filter interfaces or subsystems. Some
occur in public API names as well as internal symbols; the surrounding finding
names the relevant path.

\item[HDF5 object / attribute / link / reference.]
Groups, datasets, and committed datatypes are principal HDF5 objects. An
attribute is a named metadata item with its own datatype and dataspace, attached
to an object. A hard link gives an object a name in a group; soft and external
links instead store targets resolved during traversal and may dangle. An object
reference identifies an object, while a region reference also identifies a
selection within a dataset.

\item[Heap.]
In ordinary programming, a heap is the memory area used for dynamic
allocation. In the HDF5 format, local, global, and fractal heaps are also named
on-disk metadata structures. The surrounding phrase determines which meaning
applies.

\item[Hostile / malicious / untrusted input.]
Input for which the processor cannot rely on the originator's intentions or
care. \emph{Untrusted} includes accidentally malformed third-party files;
\emph{malicious} means intentionally crafted. Safe parsing should not depend
on knowing which one arrived.

\end{description}

\subsubsection{I--P}
\label{glossary:ip}

\begin{description}[leftmargin=1.5em,style=nextline]

\item[Immutable file identity.]
See \emph{Open handle / immutable file identity}. It is the property used to
show that validation and later processing apply to the same bytes, rather than
merely to the same path name.

\item[Impact.]
The I1--I4 axis describing the greatest credible consequence supported by the
evidence, from limited through severe or systemic. It is not an unconstrained
worst case and is not multiplied numerically by exposure.

\item[Independent copy.]
A separately stored copy whose failure modes and administrative control are not
the same as the working file. For unregenerable data, a duplicate on the same
full device or in the same failure domain is not the independent recovery
control intended by the roadmap.

\item[Independent implementation.]
A separately developed reader or writer that interprets the HDF5 format
without simply calling the reference C library for the relevant work. It can
reveal specification or reference-library errors, but it has its own parser,
validator, writer, and dependency risks.

\item[Inherent risk.]
Risk for the stated scenario and deployment before optional, proposed, or
organization-specific controls receive credit. Only mandatory controls
intrinsic to the identified implementation, version, and default configuration
are included.

\item[Input validation.]
Checking that input has an allowed type, range, size, relationship, and
structure before it influences trusted operations. Local field checks are not
enough when an invariant spans an object or reference graph.

\item[Integer overflow / divide by zero.]
Integer overflow occurs when an arithmetic result cannot be represented in the
destination integer type; division by zero has no valid arithmetic result.
File-derived dimensions, counts, and sizes must be checked before either result
can influence allocation, address calculation, or copying.

\item[Integrity.]
For data, see \emph{Data integrity}. In a software supply chain, integrity is
evidence that an artifact has not changed since an identified producer
published or signed it. Either meaning must state its object and boundary;
integrity is not the same as provenance or safety.

\item[Interoperability.]
The ability of different tools, versions, and implementations to exchange and
interpret data as intended. It depends on conformance, supported features,
extensions, configuration, and compatibility policy.

\item[Interpretability.]
The ability to recover the intended meaning of preserved data, not merely its
bytes. It can depend on schema, provenance, format bounds, external resources,
filters, connectors, drivers, and a reader that implements the relevant rules.

\item[Invariant.]
A condition that must always hold for a structure or operation to be valid,
such as a count fitting its buffer or two metadata fields agreeing. An
invariant registry ties each condition to enforcement, tests, fuzz targets,
findings, and migration status.

\item[Isolation / process isolation.]
Separating file processing from sensitive resources, other processes, or the
host so a failure has limited reach. Containers, restricted workers, operating-
system sandboxes, and separate accounts are possible mechanisms; their actual
boundaries must be verified.

\item[Journal / write-ahead log.]
A journal records update information used to recover consistent state. A
write-ahead log records the intended change before applying it to the main
data. These are candidate crash-consistency mechanisms, not guarantees HDF5 is
assumed to provide today.

\item[Least privilege.]
Giving a process only the filesystem, network, credential, code-loading, and
other authority required for its task, for only as long as needed. It limits
the consequence when an input or component misbehaves.

\item[Legacy.]
An older format, feature, file, tool, or behavior retained for compatibility.
Legacy status does not mean unsafe, but historical tolerance of invalid input
must not silently become a permanent compatibility requirement.

\item[Library / reference library / \code{libhdf5}.]
A library is reusable program code. \code{libhdf5} is the principal HDF5 C
library; \emph{reference library} emphasizes its central role in implementing
the format. Reference status does not make its behavior automatically
normative when it conflicts with the specification.

\item[Likelihood / severity in SSP SIG registries.]
Likelihood and severity are separate 1--5 inputs in the SSP SIG mechanism
registries, where they are multiplied and banded. They are not this report's
exposure and impact levels, and their product cannot be converted into this
report's scenario-based risk tiers.

\item[Linkage: static, dynamic, system, bundled, or vendored.]
Static linkage copies library code into a binary at build time; dynamic linkage
loads a shared library at run time. A system library comes from the operating
environment; bundled or vendored code is shipped within another project. The
linkage determines which HDF5 code a user actually runs and how it is updated.

\item[Logical size / physical size.]
Logical size is the amount of data described after interpreting dimensions and
types; physical size is the storage currently occupied. Compression, sparse
layout, fill values, external storage, and VDS mappings can make them differ by
orders of magnitude.

\item[Lossless / lossy compression.]
Lossless compression reconstructs the original values exactly. Lossy
compression intentionally discards information according to precision, rate,
or error settings. Those settings and the resulting sizes can reveal data
characteristics even when values are not directly shown.

\item[Machine-readable specification.]
A structured description that software can process to generate validators or
fixtures. It makes choices testable but becomes normative only through an
explicit governance decision and cannot resolve ambiguity by itself.

\item[Malformed / invalid / inconsistent file.]
A file whose bytes violate required structural, arithmetic, or relational
conditions. \emph{Malformed} does not imply malicious intent. A valid but
policy-disallowed file is different: it may conform to the format yet still be
unsafe for a particular deployment.

\item[Managed runtime / native code.]
A managed runtime, such as a JVM or .NET runtime, supplies services such as
automatic memory management. Native code executes directly under platform
conventions, commonly through C or C++. A binding can cross from managed code
into native HDF5 code and therefore inherit both models' failure modes.

\item[Manifest.]
A machine-readable list of expected artifacts, identities, versions,
capabilities, or dependencies. A protected plugin manifest can support an
allowlist; a corpus manifest can state each fixture's provenance and expected
outcome.

\item[Materialization / activation.]
Materialization constructs an in-memory object from validated data. Activation
causes effects such as publishing it to a cache, registering an identifier,
loading a plugin, decompressing data, invoking a callback, or opening an
external file. Successful parsing must not automatically authorize activation.

\item[Memory corruption / memory safety.]
Memory corruption changes or accesses memory outside the program's intended
rules. Memory safety prevents classes such as out-of-bounds access,
use-after-free, double-free, and invalid lifetime use. A memory-safe language
reduces many such defects but does not prevent resource exhaustion, logic
errors, unsafe foreign code, or policy failures.

\item[Memory leak.]
Allocated memory that is no longer usable but has not been released. Repeated
leaks can exhaust a process even when no individual access crosses a memory
boundary.

\item[Metadata / semantic metadata / structural metadata.]
Metadata describe data. Semantic metadata convey meaning, such as names,
attributes, units, provenance, or enumeration labels. Structural metadata
describe organization and storage, such as hierarchy, shapes, layouts, chunk
indexes, filters, addresses, and references. Either can be sensitive.

\item[Metadata cache / cache image.]
The metadata cache retains decoded HDF5 metadata in memory. A cache image is a
serialized representation used to accelerate a later open. Both manage
interrelated structures whose corrupted sizes, ranks, or lifetimes can affect
parser and recovery paths.

\item[Migration / repair / regeneration.]
Migration transforms data to a supported representation. Repair changes an
existing artifact to restore a defined valid state. Regeneration recreates it
from an authoritative source. Each must be deterministic, preserve provenance,
and state how checksums and meaning change.

\item[Mitigation / treatment.]
An action intended to reduce likelihood, exposure, or impact, or to avoid,
transfer, contain, or recover from a scenario. A proposed mitigation is not
treated as effective until acceptance evidence supports it.

\item[Model artifact / model loader.]
A stored machine-learning model, weights file, or related package and the
software that reconstructs it. An HDF5 container can hold application-level
instructions or objects whose security semantics are defined by the loader,
not by HDF5 alone.

\item[Named VFD families.]
SEC2 and STDIO use ordinary local-file mechanisms; Core uses memory; Family,
Split, Multi, and Subfiling distribute one logical file across physical files;
MPI-I/O supports parallel access; ROS3 reads object storage; and Mirror, Log,
and Onion add replication, logging, or history. These names identify storage
behavior, not general meanings of words such as ``core'' or ``family.''

\item[Normative / informative / implementation-defined.]
Normative text states requirements for conformance. Informative text explains
without imposing a requirement. Implementation-defined behavior is deliberately
left to an implementation but should be documented with its compatibility
consequences.

\item[Null dereference.]
Attempting to access memory through a null pointer. The usual result is a crash,
but the precise consequence depends on the code and platform.

\item[Object header / superblock / B-tree / free-space metadata.]
Internal on-disk structures used to describe an HDF5 file. The superblock
identifies global file properties; object headers hold messages describing
objects; B-trees index records; free-space metadata tracks reusable regions.
These structures are parser inputs even though most users never manipulate
them directly.

\item[Object store.]
A storage service that addresses data as named objects rather than ordinary
filesystem byte streams. S3-compatible stores are examples; their credentials,
network requests, copies, logs, and consistency behavior become part of the
deployment boundary.

\item[On-disk.]
Stored in the serialized file representation rather than existing only in
program memory. On-disk metadata remain input-controlled when a file comes from
an untrusted source.

\item[Open handle / immutable file identity.]
An open handle is a process reference to an already opened file or resource.
Binding validation and later use to the same handle, or to a cryptographically
or otherwise reliably verified immutable identity, prevents a path from being
swapped to different bytes between the two steps.

\item[Open or processing profile.]
The report proposes named bundles of policy choices and resource limits for
opening or processing a file, with design targets such as legacy,
\emph{trusted-fast}, \emph{untrusted-strict}, or forensic. These names are not
current HDF5 guarantees; a profile becomes useful only when the relevant APIs
and runtime paths enforce it.

\item[Oracle.]
A source of expected answers used in testing, such as a validator that labels a
fixture valid or invalid. An oracle can expose a disagreement but must itself
be reviewed; its answer is not automatically the governing specification.

\item[Parser.]
Code that reads bytes and constructs meaning from their structure. Because
HDF5 is a complex binary format, its library and tools contain many interacting
parsers for metadata and data descriptions.

\item[Path traversal / path confinement.]
Path traversal uses components such as \code{..} or absolute paths to escape an
intended directory. Confinement ensures resolved external resources remain
under an explicitly permitted root after normalization and resolution.

\item[Payload.]
The primary data values a user intends to store, as distinguished from
metadata, structure, indexes, logs, caches, and other context. Payload-only
protection can therefore leave substantial information exposed.

\item[Penetration test / exploit development / source review.]
A penetration test actively attempts to compromise a system. Exploit
development turns a weakness into a reliable attack technique. Source review
examines code for defects or assurance gaps. This report is a targeted
technical examination and does not claim exhaustive performance of any of
these activities.

\item[Plugin.]
Code loaded to extend a program without rebuilding it. HDF5 filter, VOL, and
some VFD extensions can be plugins. A plugin is executable code, not inert
file data, and normally runs with the host process's authority.

\item[Poisoned or suspect file state.]
A proposed durable indication that a write failure may have left a file unsafe
to trust. Clearing such a state would require a defined diagnostic and recovery
procedure rather than an ordinary reopen.

\item[Precondition.]
A fact that must hold for a scenario to occur, such as a feature being enabled,
the process having access to a target, or an affected version being deployed.
Preconditions shape exposure; they are not implementation-effort factors.

\item[Preflight validator / enforcement coupling.]
A preflight validator inspects a file before expensive allocation,
deserialization, plugin activation, or external access. Its result is advisory
unless it is bound to the same immutable file identity or open handle and the
same policy is rechecked at side-effect boundaries.

\item[Preservation object.]
Everything that must be retained to make archived information accessible and
interpretable. It may be a self-contained HDF5 file, or it may also require
external stores, codecs, connector behavior, configuration, schemas, test
vectors, credentials policy, and a reproducible reader environment.

\item[Primary artifact / primary evidence.]
The source closest to the fact being established, such as the specification,
affected source and fix, test result, vendor advisory, or first-party incident
record. A primary source can still be incomplete or wrong and must be bounded
to what it actually shows.

\item[Primary assessment object / boundary case study / contextual reference.]
The primary object is directly assessed. A boundary case study assesses only a
defined interaction with that object. A contextual reference explains the
surrounding ecosystem without making an audit conclusion about the referenced
product. See also \emph{Assessment surface / scope level}.

\item[Privacy.]
Protection against adverse consequences from disclosure, retention, inference,
aggregation, correlation, surveillance, or secondary use of information.
Privacy harm can occur during authorized, technically correct operation and is
not limited to a security breach or legal compliance question.

\item[Privilege escalation.]
Gaining authority beyond what the initiating user or process was intended to
have. Loading code through a path writable by a lower-privileged user can let
that code execute with a higher-privileged process's permissions.

\item[Profile.]
A named bundle of assumptions or settings. The report uses the word for
security/open policy, build configuration, producer/consumer compatibility,
deployment scope, and archival requirements. These profiles answer different
questions and are not interchangeable.

\item[Property list.]
An HDF5 API object that carries configuration for operations such as opening a
file, creating a dataset, selecting a driver, or loading external resources.
Several proposed security controls are expressed as properties so policy can
be bound to the relevant operation.

\item[Provenance.]
Evidence of where data or software came from, who or what produced it, which
versions and transformations were involved, and how custody changed. A valid
signature can support provenance but does not describe every transformation or
prove correctness.

\end{description}

\subsubsection{Q--Z}
\label{glossary:qz}

\begin{description}[leftmargin=1.5em,style=nextline]

\item[Rank / dimension / shape / extent.]
Rank is the number of axes in an HDF5 dataspace. Dimensions give the size along
those axes; their ordered values form the shape or current extent. These
file-declared values must be checked before a consumer allocates an array.

\item[Rating status: Rated, Provisional, NE, and N/A.]
Rated means the evidence supports the impact, exposure, and tier. Provisional
means those can be estimated but material uncertainty remains. NE means a
required rating input is insufficient, and N/A means the item is not
risk-bearing. Only Rated and Provisional records carry report-wide tiers.

\item[Reader acceptance.]
The observed fact that one reader opens or processes a file. Acceptance does
not by itself prove specification conformance, internal consistency, safe
deployment, or acceptance by another version or implementation.

\item[Reader / writer / producer / consumer.]
A writer or producer creates an HDF5 artifact; a reader or consumer interprets
it. The terms can refer to libraries, tools, applications, or independently
implemented software, and each exact version and profile matters.

\item[Recommendation.]
A proposed control outcome identified separately from a finding. Its number
does not encode risk, urgency, effort, dependency, or implementation sequence.

\item[Recovery / recoverability.]
The ability and documented procedure to return data or service to a known,
usable state after failure. Reopening a file is not sufficient evidence if it
silently returns incorrect data.

\item[Recurrence / prevalence.]
Recurrence is repeated observation of a mechanism or event. Prevalence is how
widely it exists in the relevant installed population. Several CVEs can show a
recurring class without establishing current deployment prevalence.

\item[Reference implementation.]
The implementation maintained as the principal realization of a specification.
It is important evidence and often defines practical compatibility, but its
behavior does not override a clear normative requirement without an
authoritative specification decision.

\item[Regression / regression test.]
A regression is the return of a previously corrected defect or an unintended
loss of supported behavior. A regression test is a permanent, preferably
small test intended to detect that return. Stricter rejection of invalid input
is not automatically a regression.

\item[Re-identification / secondary use.]
Re-identification connects data thought to be anonymous or de-identified to a
person or entity. Secondary use applies information for a purpose beyond the
one for which it was collected or shared. Both can create privacy harm without
breaking file encryption or access controls.

\item[Release build / debug build.]
A release build is configured for supported distribution and commonly disables
assertions or adds optimization. A debug or assertion-enabled build favors
diagnostics. Security-relevant behavior must be tested in the configurations
users actually receive, including paired assertion-enabled and
\code{-DNDEBUG} builds. A \emph{supported release} is instead a shipped version
within the project's support policy; the two phrases are not interchangeable.

\item[Reproducible build / reproducible environment.]
A reproducible build can be recreated from recorded source, tools, dependencies,
and instructions with a defined matching result. A reproducible environment
preserves enough of the runtime stack to execute or reimplement a dependency
later. Either claim must state what is reproduced and how equality is tested.

\item[Residual risk: current and target.]
Current residual risk is what remains after verified existing controls. Target
residual risk is a projection for proposed controls and is not achieved until
acceptance criteria pass. Neither replaces the inherent-risk rating.

\item[Resource budget.]
See \emph{Allocation / resource budget}. A useful budget is an enforced bound
over the relevant aggregate work or resource use, not merely a per-field size
check.

\item[Response urgency / response clock.]
The report-assigned treatment lane---Immediate, Near-term, Planned, or
Monitor---and its target horizon. Urgency considers current residual risk, time
to harm, exploitation or recurrence, and risk tolerance; dependencies and
effort may shape delivery but do not lower it. These are assessment
recommendations, not stakeholder-approved commitments until adopted and
assigned.

\item[Risk / risk rubric / risk tier.]
Risk is the possibility and consequence of a stated adverse scenario. The
rubric is this report's decision method; it combines ordinal impact and
exposure through a matrix to yield Low, Moderate, High, or Critical. A tier
belongs to a bounded scenario, not to a CVE count, recommendation, technology,
or project. A range such as Moderate--High remains a range even when response
is conservatively routed at its upper bound.

\item[Risk acceptance / risk owner / risk tolerance.]
Risk acceptance is an accountable decision to retain verified remaining risk.
The risk owner has authority to make or escalate that decision. Risk tolerance
states how much risk an organization permits; silence, an expired exception, or
an unfinished mitigation is not acceptance.

\item[Root cause / symptom / consequence.]
The root cause is the earliest evidenced failure that explains an event, such
as accepting an invalid file-derived size. A crash or buffer overflow may be a
symptom or consequence. Fixing one symptom does not prove elimination of the
root-cause class or sibling paths.

\item[Runtime check.]
A condition evaluated while the deployed program runs. A security-relevant
runtime check must be present in supported release builds and return a
controlled error when it fails.

\item[Safety.]
Operational robustness and protection of integrity, availability,
recoverability, interoperability, and long-term accessibility without assuming
a malicious actor. It does not mean regulated physical or functional safety in
this report.

\item[Sandbox.]
An enforced environment that restricts what a process can read, write, execute,
or contact. A temporary directory alone is not a sandbox; the actual
filesystem, network, credential, and process boundaries must deny access.

\item[Sanitizer.]
Instrumentation that detects selected runtime errors during tests. ASan,
UBSan, and LSan cover different classes. A sanitizer-clean corpus is useful
evidence but cannot prove absence of all defects or policy failures.

\item[Scope tag.]
A label identifying where a CVE or condition belongs, such as core parser,
official tool, plugin loading, external resource, or downstream
deserialization. It supports triage but assigns neither a rating nor priority.

\item[Search path.]
See \emph{Current working directory / search path}. The order and ownership of
locations searched for external files or loadable code can determine which
resource the process actually uses.

\item[Security.]
Protection against an adversary compromising confidentiality, integrity, or
availability, or causing a process to exercise authority contrary to the
operator's intent. A security consequence can arise from a parser defect or
from unsafe use of an intentional capability.

\item[Security advisory.]
A published notice about a security issue, often stating affected versions,
impact, and treatment. Advisory identifiers such as a CVE or GHSA help locate
records; the underlying evidence and applicability still require review.

\item[Security profile.]
An enforceable bundle of feature policy and limits appropriate to a trust
context. The proposed untrusted-HDF5 profile denies unnecessary external and
executable capabilities and bounds parsing, traversal, output, and allocation.

\item[Segmentation fault / crash.]
A segmentation fault is process termination commonly caused by an invalid
memory access. A crash is any unexpected process termination. Either is an
observable failure, not by itself proof of root cause or reliable
exploitability.

\item[Selection / hyperslab / point selection.]
A selection identifies which elements of a dataspace participate in an I/O
operation. A hyperslab describes regular blocks with offsets, strides, counts,
and block sizes; a point selection lists individual coordinates. Selections
can drive large work or allocations and therefore require bounds checking.

\item[Self-describing data.]
Data that carry structural information needed to decode them, such as
data\-types, data\-spaces, names, attributes, and relationships. HDF5 can also
carry semantic metadata, but \emph{self-describing} does not guarantee that all
context needed to interpret the scientific meaning is present. The included
metadata can itself disclose sensitive context.

\item[Sensitive information.]
Information whose disclosure, inference, retention, or use can harm people,
organizations, research, operations, or obligations. Sensitivity is contextual
and can reside in metadata, structure, access patterns, and derived results as
well as payload values.

\item[Serialization.]
Turning an in-memory value or object into bytes for storage or transfer; the
reverse operation is deserialization. See also \emph{Decode / encode} and
\emph{Deserialization}.

\item[Side effect.]
An externally observable action beyond constructing a value in memory, such as
opening another file, loading code, making a network request, writing a cache,
allocating a large buffer, or invoking a callback. Validation does not by
itself authorize a side effect.

\item[Software license.]
The legal terms governing use, modification, and distribution of software.
The report's GPLv3-or-later versus permissive BSD-style discussion is a reuse
and provenance constraint, not a security rating; legal review may be required
before incorporating code or generated artifacts.

\item[Specification / normative ambiguity / adjudication.]
A specification defines the format's requirements. A normative ambiguity
exists when those requirements permit materially competing interpretations.
Adjudication is the authoritative decision that identifies the correct rule,
the artifact to change, or an explicitly implementation-defined boundary.

\item[Stack.]
In ordinary code execution, the stack holds call frames and many local values;
deep recursion can exhaust it, and a stack buffer can be overrun. In extension
chapters, a \emph{connector stack} instead means multiple layered VOL
connectors.

\item[Static analysis / dynamic testing.]
Static analysis examines source or binaries without executing the target
program; dynamic testing runs it with selected inputs. They find overlapping
but different problems, and evidence from one should not be described as if it
came from the other.

\item[Storage exhaustion.]
Running out of space or quota while writing, flushing, or closing a file. An
operation can report failure yet still leave uncertain durable state, which is
why monitoring, headroom, error signaling, and recovery are separate controls.

\item[Storage layout.]
The principal dataset layouts are compact (raw data in the object header),
contiguous (one continuous region in the HDF5 address space), chunked
(separately indexed blocks), and virtual (mappings to source datasets).
External raw storage is a form of contiguous layout that places bytes in one or
more external files. Layout affects performance, compatibility, and security
policy.

\item[Supply chain.]
The people, source, dependencies, build systems, automation, packages,
distributors, and downstream consumers through which software is produced and
delivered. Supply-chain integrity does not by itself establish parser safety or
deployment compatibility.

\item[Supported release / affected version.]
A supported release is a shipped version still covered by the maintainer's
support policy. An affected version is one for which evidence establishes the
relevant condition. Development-tree behavior, an unspecified older version,
and a supported affected release are different claims.

\item[Threat model.]
A structured description of assets, actors, trust boundaries, capabilities,
assumptions, attack mechanisms, and consequences. It helps enumerate scenarios
but does not automatically assign this report's risk ratings.

\item[Tranche.]
A bounded, independently deliverable subset of a larger migration or assurance
program, such as one family of on-disk structures. Tranches support staged
implementation; their order does not create a risk rating.

\item[Triage.]
The initial classification and routing of a report or record: determining
scope, affected versions, duplicates, reachability, owner, and next evidence or
treatment action. Triage is not the same as remediation or risk acceptance.

\item[Trust boundary / trusted computing base.]
A trust boundary is where data or control passes between parties or components
with different authority or assumptions. The trusted computing base is the set
of hardware, software, configuration, and people whose correct behavior the
security property depends on. In-process plugins normally join it.

\item[Trusted.]
A context-dependent claim that may refer to an input source, verified
provenance, code admitted to the trusted computing base, or the proposed
\emph{trusted-fast} processing profile. It does not automatically mean
conforming, harmless, confidential, or free of defects.

\item[Type confusion.]
Using a value or memory object as though it had an incompatible type. This can
produce incorrect interpretation, invalid memory access, or corruption even
when the underlying bytes are within an allocated region.

\item[Undefined behavior.]
A C/C++ operation for which the language standard imposes no required result.
It can appear to work in one build and fail, corrupt memory, or be optimized in
surprising ways in another.

\item[Uninitialized value.]
Memory or a variable used before the program has assigned a valid value. Its
contents are indeterminate and can affect addresses, sizes, control flow, or
output.

\item[Use-after-free / double-free.]
Use-after-free accesses an object after its storage has been released.
Double-free releases the same storage twice. Both are memory-lifetime errors
that can cause crashes, corruption, or, in some conditions, code execution.

\item[User block.]
An optional area at the beginning of an HDF5 file reserved for application
content. It can contain arbitrary bytes, including scripts, XML, identifiers,
or other sensitive information not visible in an ordinary dataset listing.

\item[Valid.]
An overloaded word that must be qualified. Bytes can be syntactically
decodable, internally consistent, specification-conforming, accepted by a
particular reader, and authorized under an application's policy---or only some
of these. None of the first four alone proves safe deployment.

\item[Validation scope: record, object, and graph.]
Record-local validation checks one decoded structure. Object-wide validation
checks relationships across all records forming an object. Graph validation
checks references, cycles, aggregate limits, and relationships across objects.
A file can pass the first scope and fail a wider one.

\item[Vendored library.]
A copy of a dependency included and maintained inside another project's source
or package. Users may not realize they are running it, and system-library
updates may not replace it.

\item[Version-and-path gap.]
The report's term for missing evidence about which versions contain a defect,
which real entry paths reach it, and which versions or commits fix it. Until
that gap is closed, a CVE history cannot be treated as an installed-population
applicability list.

\item[Vulnerability.]
See \emph{Defect / weakness / vulnerability}. A bug, assertion site, feature,
or CVE mention is not enough by itself to establish a vulnerability in every
version or deployment.

\item[Whole-file / payload-only protection.]
Payload-only protection covers primary data values but may leave names,
attributes, shapes, filter settings, links, and internal structure readable.
Whole-file or equivalent all-metadata protection covers the serialized file at
the stated boundary, but still does not prevent post-decryption inference,
authorized misuse, logs, caches, key loss, or uncontrolled copies.

\item[Work package / roadmap phase.]
A planned bundle of implementation and evidence work with an owner, dependency,
effort, target, and acceptance condition. Its phase is a delivery construct and
does not assign or alter a finding's risk tier.

\item[Workflow artifact.]
An output created around the primary data, such as a log, redirected command
output, cache, temporary file, converted copy, snapshot, crash dump, or derived
dataset. These artifacts can change the security, privacy, and preservation
boundary.

\end{description}

\end{document}
